\icode{E-MI-SI-L3-261030-001}
光滑的半球形物体固定在水平地面上，球心正上方有一光滑的小滑轮，轻绳的一端系一小球，靠放在半球上的\(A\)点，	另一端绕过定滑轮后用力拉住，使小球静止，如图所示。现缓慢地拉绳，在使小球沿球面由\(A\)到\(B\)的过程中，半球对小球的支持力\(F_N\)和绳对小球的拉力\(T\)的大小变化情况是\paren[]

\begin{tikzpicture}[>=stealth,font=\scriptsize]
	\draw (1.2,0) arc(0:180:1.2);
	\draw[dashed,cyan] (0,0)--(0,3);
	\draw[fill=white] (0,3) circle(.1);
    \draw[magenta] (130:1.3)--node[left,black]{\(L\)}($(0,3)+(140:.1)$);
	\draw[->,magenta] ($(0,3)+(70:.1)$)--++(-20:.6)node[right]{\(F\)};
	\shade[ball color=gray] (130:1.3)node[left]{\(A\)} circle(.1);
	\foreach \x in {-1.8,-1.6,...,1.6}
		\draw[gray] (\x,-.2)--++(.2,.2);
	\foreach \x in {-.6,-.4,...,.6}
		\draw[gray] (\x,3.6)--++(.2,-.2);
	\draw[thick] (-1.8,0)--(1.8,0) (-.6,3.4)--(.6,3.4) (0,3)--(0,3.4);
	\fill (0,3) circle(.05) (0,0) circle(.05);
	\node at(.2,1.4){\(B\)};
\end{tikzpicture}
\begin{choices}
	\item\(F_N\)变大，\(T\)变小
	\item\(F_N\)变小，\(T\)变大
	\item\(F_N\)变小，\(T\)先变小后变大
	\item\(F_N\)不变，\(T\)变小
\end{choices}
\begin{answer}
	D\qquad 提示：如图所示，以重力\(mg\)为基础构建力的矢量三角形，与边的三角形相似，\(\frac{mg}{OC} = \frac{N}{OA} = \frac{T}{AC}\)

    \begin{tikzpicture}[>=stealth,font=\scriptsize]
	\draw[gray] (1.2,0) arc(0:180:1.2);
	\draw[dashed,cyan!50] (0,0)--(0,3);
	\draw[fill=white] (0,3) circle(.1);
    \draw[magenta!50] (130:1.3)--node[left,black]{\(L\)}($(0,3)+(140:.1)$);
	\draw[->,magenta!50] ($(0,3)+(70:.1)$)--++(-20:.6)node[right]{\(F\)};
	\shade[ball color=gray!50] (130:1.3)node[above right]{\(A\)} circle(.1);
	\foreach \x in {-1.8,-1.6,...,1.6}
		\draw[gray!50] (\x,-.2)--++(.2,.2);
	\foreach \x in {-.6,-.4,...,.6}
		\draw[gray!50] (\x,3.6)--++(.2,-.2);
	\draw[thick,gray] (-1.8,0)--(1.8,0) (-.6,3.4)--(.6,3.4) (0,3)--(0,3.4);
	\fill[gray] (0,3)node[below right,black]{\(C\)} circle(.05) (0,0) circle(.05);
	\node at(.2,1.4){\(B\)};
    \draw[dashed,teal] (0,0)node[below,black]{\(O\)}--(130:1.3);
    \draw[->] (130:1.3)--++(0,-2)node[below]{\(mg\)};
    \draw[->,teal] ($(130:1.3)+(0,-2)$)--++(130:.83)coordinate(o)node[below left]{\(N\)};
    \draw[->,magenta] (o)--(130:1.3)node[left]{\(T\)};
\end{tikzpicture}
\end{answer}